% Set spellcheck for TexStudio
%!TeX spellcheck = en_US
%!TeX root = thesis.tex

\chapter{Implementation} \label{Implementation}

Here we want to describe the developed algorithms. The idea was to create the union $U$, 
where every processor $r$ has calculated a subset $S_r\subset C$ from a known set $C$ of fixed size. The union
\begin{align}
	U = \bigcup_{r=0}^{p-1} S_r
\end{align}
should be available for every processor. $p$ is here the number of processes. \\
A straight forward approach to tackle this problem is to create a bitmap for each process, give the elements of $C$ an explicit ordering
and map each element to a bit of the bitmap. An element $e$ is then defined as being in the subset $S_r$, if the the corresponding 
bit is set to 1 and set to 0 otherwise. The union can than easily be calculated by performing a bitwise OR-operation over all 
bitmaps from all processes, as the union over all subset $S_r$ is defined as
\begin{align}
	U = \{x \in C | x \in S_0 \lor x \in S_1 \lor \cdots \lor x \in S_{p-1} \}.
\end{align}
This can be done with MPI using the provided Allreduce function. To store $N$ vertices we need $N$ bits in the bitmap.
As a word size of 64-bit integer is quite common on many architectures, we need $\ceil{\frac{N}{64}}$ 64-bit integer
to represent a subset $S_r$. This can be wasteful, if the cardinality of this subset $S_r$ is quite small and most bits of the bitmap are set to 0. We can introduce here a simple sparse format to save memory and therefore communication bandwidth. By using unsigned 32-bit integer, we can represent subsets of $C$ as an explicit list of the selected integer. Here we only can represent $2^{32}$ distinct elements so the cardinality of $C$ is bounded by $2^{32} \approx 4.2\cdot10^9$. The benefit is now, that the size of the sparse representation only depends on the cardinality of the subset $S_r$. We can calculate the memory footprint $M_\mathrm{Sparse}$ of this representation as being 
\begin{align}
	M_\mathrm{Sparse} = m \cdot \qty{32}{\bit},
\end{align}
where $m = |S_r|$ is the cardinality of $S_r$. With the approximate memory footprint of the dense representation $M_\mathrm{Dense} = N \cdot \qty{1}{bit}$, we can estimate, that the sparse variant is more memory efficient, if
\begin{align}
	m \leq \frac{1}{32} = 3.125 \%.
\end{align}
In the next sections we propose some algorithms, that take advantage of the sparse format to perform such an union operation of subsets $S_r$ of a fixed set $C$.


\section{An Algorithm for $2^k$-number of processes}
The core structure for algorithm \ref{algo:ExactHalfingAllreduce} is based on the  Bruck/Desimination algorithm described in section \ref{sec:BrucksAlgo}. When only the dense representation is used (i.e. when it is chosen by the algorithm, that this would be the more efficient way to go) the algorithms are equivalent. As in the Bruck algorithm there are two phases. In the first one we perform a distributed Reduction (that's an \texttt{MPI\_Reduce\_scatter()} in the MPI standard). This phase is followed by an All-Gather (\texttt{MPI\_Allgather()} in MPI). It is easy to verify, that an distributed Reduce combined with an All-Gather results in an All-Reduce and is sketched in section \ref{sec:BrucksAlgo}. 

\begin{algorithm}
	\caption{Exact Halfing Allreduce with sparse adaptation.\\
		$r$: rank Id of the process\\
		$p$: number of processes\\
		$I_s$: input array with sparse representation of $S_r$\\
		$sparseSize$: length of sparse list\\
		$I_d$: input array with dense representation of $S_r$\\
		$bitarraySize$: count of words for bitarray\\
		$isDense$: bool to indicate, which representation holds} \label{algo:ExactHalfingAllreduce}
	\begin{algorithmic}[1]
		\Ensure $p = 2^k$ for some $k$ \Comment{Only works for power of 2 number of ranks}
		\Procedure{VS\_ExactHalfing}{$r, p, I_s, sparseSize, I_d, bitarraySize, isDense$}
		\State critSize $\gets$ \Call{getCritSize}{$bitarraySize$}
		\Comment{Calculate Critical size}
		\\
		
		\For{(shift $\gets p/2$, shift $\geq1$, shift $\gets$ shift/2)} \Comment{Reducescatter phase}
			\State commNeigh $\gets r$ XOR shift \Comment{XOR is bitwise XOR}
			\If{isDense}
				\State startBlock $\gets$ (commNeigh/shift) $\cdot$ shift
				\State endBlock $\gets$ startBlock + shift
				\State Send($I_d$[startBlock, $\dots$, endBlock - 1], "dense") \Comment{Send block. Mark message as "dense"}
			\Else
				\State Send($I_s$[0, $\dots$, $sizeSparse$-1], "sparse") \Comment{Send whole sparse array. Mark message as sparse}
			\EndIf
			\\
			\State status $\gets$ \Call{probe}{commNeigh} \Comment{Get message type and size information}
			\If{(status.tag = "dense")} \Comment{Dense block is received}
				\State startBlock $\gets$ ($r$/shift) $\cdot$ shift
				\State endBlock $\gets$ startBlock + shift
				\If{$isDense$} \Comment{Own structure is also "dense"}
					\State Recv($\tilde{I}_d$[0, $\dots$, status.count - 1], "dense")\Comment{Receive in buffer}
					\State Wait($I_d$) \Comment{Wait until it is safe to use original array}
					\ForAll{i = startBlock, $\dots$, endBlock - 1}
						\State $I_d$[i] $\gets$ $I_d$[i] $\lor$  $\tilde{I}_d$[i - startBlock] \Comment{$\lor$ is here bitwise OR}
					\EndFor
				\Else \Comment{Own structure is "sparse"}
					\State Recv($I_d$[startBlock, $\dots$, startBlock + status.count - 1], "dense")
					\ForAll{element $\in I_s$} \Comment{Add all old sparse elements in new dense array}
					\State \Call{Bitmap\_Set}{$I_d$, element}		
					\EndFor					
					\State isDense $\gets$ true
				\EndIf
			\\
			\ElsIf{(status.tag = "sparse")} \Comment{Sparse array is received}
				\If{($isDense$)}
					\State Recv($\tilde{I}_s$ [0, $\dots$, ], "sparse")
					\ForAll{element $\in \tilde{I}_s$} \Comment{Add all new sparse elements in old dense array}
					\State \Call{Bitmap\_Set}{$I_d$, element}		
					\EndFor
				\Else
					\State Recv($I_s$[$sizeSparse$, $\dots$, $sizeSparse$+status.count-1], "sparse") \\ \Comment{Receive directly in sparse array behind the old elements}
					\State $sizeSparse \gets sizeSparse$ + status.count
					\If{($sizeSparse > critSize$)} \Comment{Transform into dense, if too big}
					\ForAll{element $\in I_s$} 
					\State \Call{Bitmap\_Set}{$I_d$, element}		
					\EndFor
					\EndIf					
				\EndIf
			\EndIf			
		\EndFor
		\\
		\If{($isDense$ = false)} \Comment{Return early, if variant is still sparse}
			\State \Return
		\EndIf
		\\
		\Call{Allgather}{$I_d$}
		\EndProcedure
	\end{algorithmic}
\end{algorithm}



